\section{Results}

\subsection{Execution success}

Fully pinned projects completed successfully in 96.7\% of perturbed rebuilds.
Success decreased to 82.4\% under partial specification and 61.3\% without
version constraints (Table~\ref{tab:results}). The largest class of failures in
the unpinned condition was an incompatible transitive dependency; system-library
mismatches were the next most common. Figure~\ref{fig:success} shows that the
ranking remained stable as perturbation intensity increased.

\begin{table}[tbp]
  \centering
  \caption{Rebuild outcomes by environment-documentation strategy. Intervals
  are project-clustered 95\% bootstrap intervals. Numerical deviation is
  summarized only among successful rebuilds.}
  \label{tab:results}
  \begin{singlespace}
  \small
  \begin{tabularx}{\textwidth}{>{\raggedright\arraybackslash}X
      S[table-format=2.1]
      c
      S[table-format=2.1]
      S[table-format=2.1]}
    \toprule
    {Strategy} & {Success (\%)} & {95\% interval} & {Median deviation (\%)} & {Reproduced (\%)} \\
    \midrule
    Fully pinned       & 96.7 & 95.8--97.5 & 0.8  & 92.1 \\
    Partially specified& 82.4 & 80.6--84.0 & 4.7  & 70.6 \\
    Unpinned           & 61.3 & 58.9--63.7 & 12.9 & 43.8 \\
    \bottomrule
  \end{tabularx}
  \end{singlespace}
\end{table}

\begin{figure}[tbp]
  \centering
  \begin{tikzpicture}
    \begin{axis}[
      width=0.88\textwidth,
      height=7.2cm,
      xmin=1,xmax=5,
      ymin=35,ymax=100,
      xtick={1,2,3,4,5},
      xlabel={Perturbation intensity},
      ylabel={Successful rebuilds (\%)},
      grid=major,
      grid style={gray!20},
      legend style={draw=none,fill=none,font=\small,at={(0.02,0.02)},anchor=south west},
      tick label style={font=\small},
      label style={font=\small}
    ]
      \addplot+[color=tandfblue,mark=*,very thick]
        coordinates {(1,99.2) (2,98.4) (3,97.1) (4,95.5) (5,93.3)};
      \addlegendentry{Fully pinned}
      \addplot+[color=tandfteal,mark=square*,very thick]
        coordinates {(1,94.0) (2,89.1) (3,83.4) (4,76.7) (5,68.8)};
      \addlegendentry{Partially specified}
      \addplot+[color=tandfgold,mark=triangle*,very thick]
        coordinates {(1,82.6) (2,72.4) (3,61.8) (4,50.2) (5,39.5)};
      \addlegendentry{Unpinned}
    \end{axis}
  \end{tikzpicture}
  \caption{Execution success across increasing environmental perturbation.
  Points are simulation means; connecting lines aid visual comparison.}
  \label{fig:success}
\end{figure}

\subsection{Complementary safeguards}

Environment capture and seed control affected different outcomes. Pinning
dependencies increased the probability that a pipeline completed, whereas seed
control primarily reduced dispersion among pipelines that did complete. In the
fully pinned condition, recording a seed reduced median numerical deviation from
1.3\% to 0.4\%. In the unpinned condition, the corresponding reduction was from
14.1\% to 11.8\%, because dependency changes dominated the remaining variation.
The interaction estimate was small and its interval included zero, supporting a
complementary rather than substitutive interpretation.
